\section{Statement, inputs, and antecedents}\label{sec:context}
Let $\Sigma_g$ be a closed connected oriented surface of genus $g$. Put
\[
 \R=\Hom(\pi_1(\Sigma_g),\SL_r(\C)),\qquad
 \X=\R\mathbin{/\!\!/}\SL_r(\C),
\]
where the latter is the affine character variety. The group $\Mod(\Sigma_g)$
consists of orientation-preserving mapping classes and acts on $\X$ by
precomposition; inner automorphisms of the surface group have no effect on
characters. We use $[x,y]=xyx^{-1}y^{-1}$ throughout.

\begin{theorem}\label{thm:main}
For every pair of integers $g\geq2$ and $r\geq2$, there exist a number field
$L\subset\C$ and a homomorphism
\[
 \rho:\pi_1(\Sigma_g)\longrightarrow\SL_r(\OO_L)
\]
such that $\Mod(\Sigma_g)\cdot[\rho]$ is Zariski dense in $\X$.
\end{theorem}

Here $\OO_L$ is the full ring of integers. The conclusion concerns one
orbit of an actual integral representation and the entire complex character
variety. There is no generic-input condition and no restriction to a marked
locus. No assertion for punctured surfaces is made. For $r\geq3$ the field
in our construction depends only on $r$, although the selected matrices may
depend on $g$.

\subsection{The credited rank-two case}
Coccia--Litt \cite[Propositions 3.1.10 and 5.0.1, Theorem 5.0.2]{CL}
already prove the assertion for $r=2$, in a more general relative setting.
Their Proposition 5.0.1 produces an actual integral representation that is
good for a pants decomposition, and Proposition 3.1.10 proves density of its
mapping-class orbit. To meet the unit hypothesis one can use
$K=\Q(\sqrt5)$ and $v=(1+\sqrt5)/2$: both $v$ and $v-v^{-1}=1$ are units,
and $v$ is not a root of unity. Their field extension and construction
therefore apply to every closed genus $g\geq2$. Density over $\overline\Q$
remains density after extending scalars to $\C$: for each finite-dimensional
space of polynomials of bounded degree, the evaluation conditions have
coefficients in $\overline\Q$, so their kernel commutes with extension of
scalars. The remainder of the proof assumes $r\geq3$.

Litt's survey \cite[Theorem 5.4.13]{Litt} gives another rank-two
single-orbit route, using a Shimura-curve local system with a unitary Galois
conjugate and Previte--Xia's Theorem 1.4 \cite{PreviteXia}.
Golsefidy--Tamam \cite[Section 3.3 and Theorems G--J]{GT}
identify exceptional loci that must be excluded in the general formulation
of that Previte--Xia theorem. We record both the antecedent and this caveat;
we do not use the uncaveated statement as an input to our rank-two proof.
Golsefidy--Tamam's Theorem I gives a substantially stronger rank-two
orbit-closure theorem, under its sufficiently-dense, nonexceptional,
infinite-orbit hypotheses. Their detailed orbit-closure preprint
\cite{GTpre} explicitly treats $\SL_2$.

\subsection{What the higher-rank argument adds}
The pinned manuscript \cite{Source027} proves potential density of the
\emph{whole set} of integral characters in every rank, already over one
number field. Its closed-surface argument first puts a whole fixed-spectrum
fiber in that closure, then uses twist-torus saturation, Baire avoidance,
and a submersion. Those operations are antecedents here; density of a union
of arithmetic orbits does not itself choose one orbit.

Goldman \cite{Goldman} supplies the symplectic and Hamiltonian twist
machinery. Pickrell--Xia \cite[Proposition 2.1.24 and Lemma 2.1.34]{PX}
provide a close antecedent to infinitesimal spanning: centralizers of
products $hg^n$, dense torus powers, and a BCH/root-component argument,
followed by transverse sewing in their Section 3. Their result concerns
compact groups and generic fibers in an almost-everywhere argument.
The task here is to certify a particular complex family with two commuting
handle matrices inside the closure of one arithmetic orbit.

For comparison, the commuting twist tori associated to a pants decomposition
have dimension at most $(3g-3)(r-1)$. For $r\geq3$,
\[
 (3g-3)(r-1)<(g-1)(r^2-1)=\tfrac12\dim\X.
\]
Thus the rank-two pants-torus fiber argument does not simply give the required
higher-rank cotangent certificate. We use two noncommuting families of simple
curves and reduce their common kernel to
\begin{equation}\label{eq:residual-intro}
 (\Ad_A-\Ad_C)S=0,\qquad (I-\Ad_C)S\in\h+\hc.
\end{equation}
A symmetric/skew-symmetric degeneration proves this system generically
injective. Arithmetic parameters are then chosen in finitely many nonempty
opens. A further complex parameter, already inside the arithmetic orbit
closure, is chosen outside countably many resonance sets.

Golsefidy--Tamam \cite[Section 2.4, second remark]{GT} explicitly suggest
extending their results to $\SL_d$ and other Chevalley groups using
Larsen--Pink. That extension is not proved there. Our existence theorem
is a special case of the anticipated broader higher-rank picture, rather
than a claim that such a picture was unanticipated. The construction and
proof are new to this project; bounded searches found no higher-rank
single-orbit integral theorem, but do not establish worldwide novelty.

The choice of point cannot be replaced by an unexamined high-rank
dense-image assertion. A MathOverflow comment by M.~Kapovich reports
non-dense mapping-class orbits of Zariski-dense high-rank unitary
representations \cite{MO}. We cite that as a discussion pointer, not as a
proved theorem used here. Independently, the actual paper of
Biswas--Koberda--Mj--Santharoubane
\cite[Proposition 1.2 and Section 5]{BKMS} constructs infinite-image
representations whose \emph{$\GL_n$ characters} are fixed by the entire
mapping class group, using TQFT representations and point pushing.
That is the precise published statement we use for comparison; it is
not the same as the stronger dense-image observation in the comment.
Neither observation is needed as a lemma in our proof.

\subsection{Geometric input and conventions}\label{subsec:geometry}
We use the following exact input from \cite{Source027}: the determinant-one
fiber
\[
 \{(A_1,B_1,\ldots,A_g,B_g)\in\SL_r^{2g}:
       \prod_{i=1}^g[A_i,B_i]=I\}
\]
is irreducible of dimension $(2g-1)(r^2-1)$. This is the
$\SL_r$ intermediate assertion established \emph{inside the proof} of
Proposition \texttt{prop:surface-irreducibility}, in
\texttt{surfaces.tex}, before the scalar-extension paragraph. The
proposition's stated conclusion concerns the $\GL_r$ fiber; we do not
infer determinant-one irreducibility from that conclusion. The continuous
surjective quotient morphism then makes $\X$ irreducible.

At a smooth irreducible representation the tangent space of $\X$ is the
kernel of the linearized surface relation modulo infinitesimal simultaneous
conjugation. The trace pairing on $\slr$ and cup product give the
nondegenerate holomorphic Goldman form on this smooth locus
\cite[Introduction and Sections 2--3]{Goldman}. We will check smoothness
at our chosen representations directly. All flow arguments take place in
an ambient neighborhood in this smooth irreducible locus, hence in
Goldman's $D/G$; the orbit closure itself need not be smooth.

Write the first two pairs of surface generators as $(a,b),(c,d_0)$ and the
others as $(a_i,b_i)$ for $3\leq i\leq g$. Their matrices are
$(A,B,C,D),(A_i,B_i)$. Let $T\subset\SL_r$ be the diagonal torus and
$\h\subset\slr$ its Lie algebra. Given a regular semisimple $C=QC_\mu Q^{-1}$,
write $\TC=QTQ^{-1}$ and $\hc=Q\h Q^{-1}$.
